% coalition.tex -- Section 7: Coalition Formation

\section{Coalition Formation}
\label{sec:coalition}

Coalition formation extends N-player games with a structured negotiation phase,
enabling agents to form binding or non-binding agreements before acting. This
tests whether agents can build trust, coordinate collectively, and handle
betrayal---capabilities directly relevant to multi-agent alignment.

% TODO: Motivate coalition formation as alignment testbed

\subsection{Two-Phase Protocol}
\label{sec:coalition-protocol}

Each round of a coalition game proceeds in two phases:

\begin{enumerate}[nosep]
    \item \textbf{Negotiate Phase}: players exchange coalition proposals,
          specifying target members, proposed joint strategies, and benefit-sharing
          arrangements. Multiple rounds of proposal and counter-proposal are
          supported.
    \item \textbf{Act Phase}: players choose actions. Coalition members may
          honor or violate their agreements, subject to the enforcement mode.
\end{enumerate}

% TODO: Protocol diagram, message format specification

\subsection{Enforcement Modes}
\label{sec:enforcement}

Three enforcement modes model different institutional settings:

\begin{description}[nosep]
    \item[Cheap Talk.] Agreements are non-binding. Players may freely
        deviate from coalition commitments with no direct penalty.
    \item[Penalty.] Defecting from a coalition incurs a configurable
        payoff penalty, modeling reputation costs or contractual damages.
    \item[Binding.] Coalition agreements are fully enforced. Once a player
        joins, their action is locked to the agreed strategy.
\end{description}

% TODO: Formal specification of enforcement mechanics

\subsection{Side Payments and Defection Detection}
\label{sec:side-payments}

The coalition framework supports side payments---transfers between coalition
members to incentivize participation. Defection detection mechanisms track
whether coalition members honored their commitments, enabling retaliatory
strategies in future rounds.

% TODO: Side payment algebra, defection detection algorithm

\subsection{Coalition Strategies}
\label{sec:coalition-strategies}

Built-in coalition strategies model different agent dispositions:

\begin{description}[nosep]
    \item[Loyal.] Always honors coalition agreements; never defects.
    \item[Betrayer.] Joins coalitions to gain trust, then defects at the
        optimal moment for maximum exploitation.
    \item[Conditional.] Honors agreements as long as other members do;
        defects in response to observed defection.
    \item[Random.] Randomly joins or rejects coalition proposals and
        randomly honors or violates agreements.
\end{description}

% TODO: Strategy pseudocode

\subsection{Coalition Game Library}
\label{sec:coalition-games}

Seven games are specifically designed for coalition dynamics:

\begin{itemize}[nosep]
    \item \textbf{Cartel Formation} --- firms collude on pricing vs.\ compete
    \item \textbf{Alliance Game} --- mutual defense pacts with defection risk
    \item \textbf{Voting Blocs} --- coordinated voting for collective benefit
    \item \textbf{Ostracism Game} --- coalition excludes non-cooperators
    \item \textbf{Resource Trading} --- bilateral exchange within coalitions
    \item \textbf{Rule Voting} --- coalitions propose and vote on rule changes
    \item \textbf{Commons Management} --- coalition-based resource governance
\end{itemize}

% TODO: Formal game specifications for each coalition game
